\documentclass{article}

\begin{document}

% Rule for a Polynomial:
A \textbf{polynomial} is an algebraic expression that consists of terms where variables are raised to non-negative integer powers and combined using addition, subtraction, and multiplication. Importantly, it does \textbf{not} involve division by any variables.

For example:

\[
P(x) = 3x^2 + 2x + 1
\]

is a polynomial, since it only involves terms with variables raised to non-negative integer powers and no division by variables.

% Rule for a Rational Expression:
A \textbf{rational expression} is an algebraic expression that can involve division by a polynomial. The \textbf{numerator} and \textbf{denominator} of a rational expression may both be polynomials, but the overall expression involves a division by a variable.

For example:

\[
R(x) = \frac{3x^2 + 2x + 1}{x + 1}
\]

is a rational expression because it involves division by \(x + 1\), which is a polynomial. However, this expression is \textbf{not} a polynomial due to the division by a variable in the denominator.

\end{document}
